\documentclass[12pt,a4paper]{article}
\usepackage[T2A]{fontenc}
\usepackage[utf8]{inputenc}
\usepackage[russian]{babel}
\usepackage{float}
\usepackage{amsmath, amssymb}
\usepackage{geometry}
\usepackage{tikz}
\usepackage{tikz-3dplot}
\usepackage{caption}
\usetikzlibrary{arrows.meta, calc, decorations.pathmorphing}
\geometry{margin=2cm}

\begin{document}

\setcounter{page}{93}

\subsection*{\textbf{ЛЕКЦИЯ 13. МНОГОЧЛЕНЫ}}

\subsection*{\textbf{13.1. Основные понятия. Деление многочленов. Теорема Безу}}

\textbf{Определение.} Многочленом $n$-ой степени называется функция
$$
P(x)=a_0x^n+a_1x^{n-1}+\ldots+a_{n-1}x+a_n,
$$
где $x\in\mathbb C$; $a_i\in\mathbb C$; $i=0,1,2,\ldots,n$; $a_0\ne0$.

\textbf{Обозначение.} $\deg P$ --- степень многочлена $P(x)$.

\textbf{Замечание.} Сложение и умножение многочленов определяются
обычным образом, при этом:
$$
\deg(P+Q)\le \max(\deg P;\deg Q);
$$
$$
\deg(P\cdot Q)=\deg P+\deg Q.
$$

\textbf{Деление многочленов}

\textbf{Теорема.} Для любых двух многочленов $f(x)$ и $g(x)$, где
$g(x)\ne0$, существует, и притом единственная, пара многочленов
$q(x)$ и $r(x)$ такая, что
\begin{equation}
  f(x)=g(x)\cdot q(x)+r(x) \tag{13.1}
\end{equation}
где либо $r(x)=0$, либо $\deg r<\deg g$.

$f(x)$ --- делимое, $g(x)$ --- делитель, $q(x)$ --- частное, $r(x)$ --- остаток.

\textbf{Доказательство}

Если $f(x)=0$, то берем $r(x)=q(x)=0$ и теорема верна.

Пусть $f(x)\ne0$.

I случай $\deg f<\deg g\Rightarrow f=g\cdot0+f$ и $\deg f<\deg g$.

II случай $\deg f\ge\deg g$
$$
f(x)=a_0x^n+a_1x^{n-1}+\ldots+a_{n-1}x+a_n;
$$
$$
g(x)=b_0x^m+b_1x^{m-1}+\ldots+b_{m-1}x+b_m.
$$
Применим метод математической индукции по степеням $n$ многочлена $f(x)$.

1) $n=0\Rightarrow m=0$ и $q(x)=\frac{a_n}{b_m}$; $r(x)=0\Rightarrow$
Теорема верна.

2) Считаем, что теорема верна для любого многочлена степени $<n$.

Обозначим
\begin{equation}
f_1(x)=f(x)-\frac{a_0}{b_0}x^{n-m}\cdot g(x) \tag{13.2}
\end{equation}
Либо $f_1(x)=0$, либо $\deg f_1<n$.

Если $f_1(x)=0\Rightarrow f(x)=\frac{a_0}{b_0}x^{n-m}\cdot g(x)$, т.е.
$$
q(x)=\frac{a_0}{b_0}x^{n-m} \text{ и } r(x)=0\Rightarrow \text{ Теорема верна.}
$$
Если $\deg f_1<n$, то по предположению индукции
$$
f_1(x)=g(x)\cdot q_1(x)+r(x),
$$
где либо $r(x)=0$, либо $\deg r<\deg g$.

Из (13.2) получим:
$$
f(x)=f_1(x)+\frac{a_0}{b_0}x^{n-m}\cdot g(x)=g(x)\cdot q_1(x)+r(x)+
$$
$$
+\frac{a_0}{b_0}x^{n-m}\cdot
g(x)=g(x)\cdot\left(\frac{a_0}{b_0}x^{n-m}+q_1(x)\right)+r(x)=
$$
$$
=\left[\frac{a_0}{b_0}x^{n-m}+q_1(x)=q(x)\right]=g(x)\cdot q(x)+r(x)\Rightarrow
$$
$$
\Rightarrow \text{ Теорема верна.}
$$
Докажем единственность разложения (13.1). Пусть существуют два разложения:

(13.1) и
\begin{equation}
f(x)=g(x)\cdot q_1(x)+r_1(x) \tag{13.3}
\end{equation}
$$
\deg r_1<\deg g.
$$
Вычитаем (13.3) из (13.1):
$$
0=g(x)(q(x)-q_1(x))+r(x)-r_1(x)\Rightarrow
$$
$$
\Rightarrow g(q-q_1)=r_1-r.
$$
Пусть $q\ne q_1\Rightarrow \deg(g\cdot(q-q_1))\ge\deg g$.

Но $\deg(r_1-r)\le\max(\deg r_1,\deg r)<\deg g$ --- противоречие
$$
\Rightarrow q_1=q \text{ и } r_1=r.
$$
\textbf{ч.т.д.}

\textbf{Пример}

$3x^4+x^3-1=f(x)$; $x^2+x+1=g(x)$.

$$
\begin{array}{rrrr|l}
3x^4&+x^3&&-1&x^2+x+1\\[-2mm]
\cline{5-5}
3x^4&+3x^3&+3x^2&&3x^2-2x-1\\
\cline{1-3}
&-2x^3&-3x^2&&-1\\
&-2x^3&-2x^2&-2x&\\
\cline{2-4}
&&-x^2&+2x&-1\\
&&-x^2&-x&-1\\
\cline{3-5}
&&&3x&
\end{array}
$$
$$
3x^2-2x-1=q(x);\qquad 3x=r(x);
$$
$$
3x^4+x^3-1=(x^2+x+1)(3x^2-2x-1)+3x.
$$

\textbf{Теорема Безу}

Остаток от деления многочлена $f(x)$ на $x-a$ равен $f(a)$.

\textbf{Доказательство}

Разделим $f(x)$ на $x-a$:
$$
f(x)=(x-a)q(x)+r(x),
$$
где $\deg r<\deg(x-a)=1\Rightarrow r(x)=r=\operatorname{const}$.

При $x=a$ получим $r=f(a)$.

\textbf{ч.т.д.}

\subsection*{\textbf{13.2. Корни многочленов. Кратность корня. Теорема Виета}}

\textbf{Определение.} $x=\alpha$ называется корнем многочлена $f(x)$,
если $f(\alpha)=0$.

\textbf{Замечание.} Если остаток от деления $f(x)$ на $g(x)$ равен 0,
то говорят, что $f(x)$ делится на $g(x)$.

\textbf{Следствие из теоремы Безу.} Если $x=\alpha$ --- корень
$f(x)$, то $f(x)$ делится на $x-\alpha$.

\textbf{Кратность корня}

\textbf{Определение.} Корень $x=\alpha$ называется корнем кратности
$k$ многочлена $f(x)$, если $f(x)$ делится на $(x-\alpha)^k$ и не
делится на $(x-\alpha)^{k+1}$.

\textbf{Теорема.} $x=\alpha$ --- корень кратности $k$ многочлена
$f(x)$ тогда и только тогда, когда
$$
f(x)=(x-\alpha)^k\cdot g(x),
$$
где $g(\alpha)\ne0$.

\textbf{Основная теорема алгебры.} Всякий многочлен, степень которого
не меньше единицы, имеет хотя бы один корень, в общем случае комплексный.

\textbf{Следствие.} Всякий многочлен степени $n\ge1$ имеет ровно $n$
корней, при этом
\begin{equation}
\begin{aligned}
P_n(x)&=a_0x^n+a_1x^{n-1}+\cdots+a_{n-1}x+a_n=\\
&=a_0(x-\alpha_1)(x-\alpha_2)\cdot\ldots\cdot(x-\alpha_n),
\end{aligned}
\tag{13.4}
\end{equation}
где $\alpha_1,\ldots,\alpha_n$ --- корни многочлена $P_n(x)$.

\textbf{Доказательство}

По основной теореме алгебры $\exists\alpha_1\in\mathbb C:
P_n(\alpha_1)=0$. Тогда по следствию из теоремы Безу
$P_n(x)=(x-\alpha_1)P_{n-1}(x)$.

Если $n-1\ge1$, то $\exists\alpha_2\in\mathbb C: P_{n-1}(\alpha_2)=0\Rightarrow$
$$
\Rightarrow P_n(x)=(x-\alpha_1)(x-\alpha_2)P_{n-2}(x) \text{ и т.д.}
$$
В результате получим разложение (13.4).

\textbf{ч.т.д.}

\textbf{Замечание.} Сгруппировав в равенстве (13.4) множители с
одинаковыми корнями, получим:
\begin{equation}
P_n(x)=a_0(x-\alpha_1)^{k_1}\cdot(x-\alpha_2)^{k_2}\cdot\ldots\cdot(x-\alpha_r)^{k_r},
\tag{13.5}
\end{equation}
$$
\alpha_i\ne\alpha_j \text{ при } i\ne j;\qquad k_1+k_2+\cdots+k_r=n;
$$
$$
k_i \text{ --- кратность корня } \alpha_i.
$$
(13.5) --- каноническое разложение многочлена (на множестве комплексных чисел).

\textbf{Теорема Виета}

Если $\alpha_1,\alpha_2,\ldots,\alpha_n$ --- корни многочлена
$$
P_n(x)=a_0x^n+a_1x^{n-1}+\ldots+a_{n-1}x+a_n,
$$
то
$$
\begin{cases}
\alpha_1+\alpha_2+\ldots+\alpha_n=-\dfrac{a_1}{a_0};\\[1mm]
\alpha_1\alpha_2+\ldots+\alpha_1\alpha_n+\alpha_2\alpha_3+\ldots+\alpha_{n-1}\alpha_n=\dfrac{a_2}{a_0};\\[1mm]
\ldots\\
\alpha_1\alpha_2\cdot\ldots\cdot\alpha_k+\ldots=(-1)^k\dfrac{a_k}{a_0};\\[1mm]
\ldots\\
\alpha_1\alpha_2\cdot\ldots\cdot\alpha_n=(-1)^n\dfrac{a_n}{a_0}.
\end{cases}
$$

\subsection*{\textbf{13.3. Многочлены с действительными коэффициентами}}

Пусть
$$
P_n(x)=a_0x^n+a_1x^{n-1}+\ldots+a_{n-1}x+a_n,
$$
$a_i\in\mathbb R$; $i=0,1,\ldots,n$.

\textbf{Теорема.} Если $x=\alpha\in\mathbb C$ --- корень многочлена
$P_n(x)$, то $x=\bar\alpha$ --- тоже корень.

\textbf{Доказательство}

Если $\alpha$ --- корень $P_n(x)$, то
$$
a_0\alpha^n+a_1\alpha^{n-1}+\ldots+a_{n-1}\alpha+a_n=P_n(\alpha)=0;
$$
$$
\overline{P_n(\alpha)}=\bar0;
$$
$$
\overline{a_0\alpha^n+a_1\alpha^{n-1}+\ldots+a_{n-1}\alpha+a_n}=0;
$$
$$
\bar a_0\bar\alpha^n+\bar a_1\bar\alpha^{n-1}+\ldots+\bar
a_{n-1}\bar\alpha+\bar a_n=0;
$$
$$
a_0\bar\alpha^n+a_1\bar\alpha^{n-1}+\ldots+a_{n-1}\bar\alpha+a_n=0;
$$
$$
P_n(\bar\alpha)=0\Rightarrow \bar\alpha \text{ --- корень } P_n(x).
$$
\textbf{ч.т.д.}

\textbf{Замечание.} Если $\alpha$ --- корень многочлена $P_n(x)$, то
$P_n(x)$ делится на $(x-\alpha)$.

$\bar\alpha$ --- тоже корень $\Rightarrow P_n(x)$ делится на $(x-\bar\alpha)$.

Таким образом, $P_n(x)$ делится на
$$
(x-\alpha)(x-\bar\alpha)=x^2-(\alpha+\bar\alpha)x+|\alpha|^2
$$
--- квадратный трехчлен с действительными коэффициентами.

\textbf{Замечание.} Если
$P_n(x)=a_0x^n+a_1x^{n-1}+\ldots+a_{n-1}x+a_n$ --- многочлен с
действительными коэффициентами, то
\begin{equation}
\begin{aligned}
P_n(x)&=a_0(x-\alpha_1)^{k_1}\cdot(x-\alpha_2)^{k_2}\cdot\ldots\cdot(x-\alpha_s)^{k_s}\cdot\\
&\quad\cdot(x^2+p_1x+q_1)^{r_1}\cdot\ldots\cdot(x^2+p_lx+q_l)^{r_l},
\end{aligned}
\tag{13.6}
\end{equation}
где $\alpha_i\in\mathbb R$; $k_i\in\mathbb N$; $i=1,\ldots,s$;

$p_j,q_j\in\mathbb R$; $r_j\in\mathbb N$; $j=1,\ldots,l$.

(13.6) --- каноническое разложение многочлена (на множестве
действительных чисел).

\textbf{Пример}

Решить уравнение
$$
6x^4+5x^3-9x^2+8x-4=0.
$$

\textbf{Решение.} Если уравнение имеет целые корни, то они содержатся
среди чисел: $\pm1$; $\pm2$; $\pm4$ --- делители свободного члена.
Проверкой убеждаемся, что $x_1=-2$.

Делим многочлен $P(x)=6x^4+5x^3-9x^2+8x-4$ на $x+2$ по схеме Горнера.
$$
\begin{array}{c|ccccc}
&6&5&-9&8&-4\\
-2&6&-7&5&-2&0
\end{array}
$$
$$
\Rightarrow P(x)=6x^4+5x^3-9x^2+8x-4=(x+2)(6x^3-7x^2+5x-2);
$$
Далее решаем уравнение
$$
6x^3-7x^2+5x-2=0.
$$
Это уравнение целых корней не имеет, но оно может иметь рациональные
корни, содержащиеся среди чисел: $\pm\frac12$; $\pm\frac13$;
$\pm\frac16$; $\pm\frac23$.

Проверкой убеждаемся, что $x_2=\frac23$.

Делим многочлен $Q(x)=6x^3-7x^2+5x-2$ на $3x-2$ «уголком».
$$
\begin{array}{rrrr|l}
6x^3&-7x^2&+5x&-2&3x-2\\[-2mm]
\cline{5-5}
6x^3&-4x^2&&&2x^2-x+1\\
\cline{1-2}
&-3x^2&+5x&-2&\\
&-3x^2&+2x&&\\
\cline{2-3}
&&3x&-2&\\
&&3x&-2&\\
\cline{3-4}
&&&0&
\end{array}
$$
$$
\Rightarrow 6x^3-7x^2+5x-2=(3x-2)(2x^2-x+1)
$$
Окончательно решаем уравнение
$$
2x^2-x+1=0
$$
$$
x_{3,4}=\frac{1\pm\sqrt{1-8}}{4}=\frac{1\pm i\sqrt7}{4}.
$$
\textbf{Ответ:} $\frac23$, $-2$; $\frac{1\pm i\sqrt7}{4}$.

\textbf{Пример}

Найти корни многочлена $p(z)=2z^4+11z^3+20z^2+7z-10$ и определить их
кратность. Разложить многочлен $p(z)$ на множители на множестве
действительных чисел и на множестве комплексных чисел.

\textbf{Решение}

$z_1=-2$:
$$
\begin{array}{c|ccccc}
&2&11&20&7&-10\\
-2&2&7&6&-5&0
\end{array}
$$
$$
\Rightarrow p(z)=(z+2)(2z^3+7z^2+6z-5);
$$
$z_2=\frac12$:
$$
\begin{array}{c|cccc}
&2&7&6&-5\\
\frac12&2&8&10&0
\end{array}
$$
$$
\Rightarrow p(z)=(z+2)\left(z-\frac12\right)(2z^2+8z+10)=
$$
$$
=(z+2)(2z-1)(z^2+4z+5) \text{ --- разложение на множители на
множестве } \mathbb R.
$$
$$
z^2+4z+5=0;
$$
$$
z=-2\pm\sqrt{4-5}=-2\pm i
$$
$$
\Rightarrow p(z)=(z+2)(2z-1)(z+2-i)(z+2+i) \text{ --- разложение на
множители на множестве } \mathbb C.
$$
$$
z_1=-2;\qquad z_2=\frac12;\qquad z_{3,4}=-2\pm i \text{ --- корни кратности } 1.
$$

\textbf{Пример}

Пусть $z_1=\frac{-1+i\sqrt{11}}{6}$ --- корень многочлена
$$
p(z)=6z^4-z^3+10z^2+2z+3.
$$
Найти остальные корни и их кратность. Разложить многочлен $p(z)$ на
множители на множестве действительных чисел.

\textbf{Решение}

1) $z_2=\frac{-1-i\sqrt{11}}{6}$
$$
(z-z_1)(z-z_2)=\left(z+\frac16-i\frac{\sqrt{11}}6\right)\cdot\left(z+\frac16+i\frac{\sqrt{11}}6\right)=
$$
$$
=\left(z+\frac16\right)^2+\frac{11}{36}=z^2+\frac{z}{3}+\frac{1}{36}+\frac{11}{36}=z^2+\frac{z}{3}+\frac13.
$$
2)
$$
\begin{array}{rrrrr|l}
6z^4&-z^3&+10z^2&+2z&+3&3z^2+z+1\\[-2mm]
\cline{6-6}
6z^4&+2z^3&+2z^2&&&2z^2-z+3\\
\cline{1-3}
&-3z^3&+8z^2&+2z&&\\
&-3z^3&-z^2&-z&&\\
\cline{2-4}
&&9z^2&+3z&+3&\\
&&9z^2&+3z&+3&\\
\cline{3-5}
&&&&0&
\end{array}
$$
$$
2z^2-z+3=0;
$$
$$
z_{3,4}=\frac{1\pm\sqrt{1-24}}{4}=\frac{1\pm i\sqrt{23}}{4}.
$$
3) Кратность у всех корней равна 1.
$$
p(z)=(3z^2+z+1)(2z^2-z+3) \text{ --- разложение на множители на
множестве } \mathbb R.
$$
$$
p(z)=6\left(z-\frac{-1-i\sqrt{11}}{6}\right)\left(z-\frac{-1+i\sqrt{11}}{6}\right)\left(z-\frac{1-i\sqrt{23}}{4}\right)\cdot
$$
$$
\cdot\left(z-\frac{1+i\sqrt{23}}{4}\right) \text{ --- разложение на
множители на множестве } \mathbb C.
$$

\textbf{Пример}

Найти корни многочлена
$$
p(z)=z^8-2z^7+z^6-2z^5+4z^4-2z^3+z^2-2z+1
$$
и их кратность. Разложить многочлен $p(z)$ на множители на множестве
действительных чисел и на множестве комплексных чисел.

\textbf{Решение}

$z_1=1$:
$$
\begin{array}{c|ccccccccc}
&1&-2&1&-2&4&-2&1&-2&1\\
1&1&-1&0&-2&2&0&1&-1&0
\end{array}
$$
$$
\Rightarrow p(z)=(z-1)(z^7-z^6-2z^4+2z^3+z-1)
$$
$z_2=1$:
$$
\begin{array}{c|cccccccc}
&1&-1&0&-2&2&0&1&-1\\
1&1&0&0&-2&0&0&1&0
\end{array}
$$
$$
\Rightarrow p(z)=(z-1)^2(z^6-2z^3+1)=
$$
$$
=(z-1)^2(z^3-1)^2=
$$
$$
=(z-1)^2(z-1)^2(z^2+z+1)^2=
$$
$$
=(z-1)^4(z^2+z+1)^2 \text{ --- разложение на множители на множестве } \mathbb R.
$$
$$
z^2+z+1=0;
$$
$$
z=\frac{-1\pm\sqrt{1-4}}{2}=\frac{-1\pm i\sqrt3}{2};
$$
$$
p(z)=(z-1)^4\left(z-\frac{-1+i\sqrt3}{2}\right)^2\left(z-\frac{-1-i\sqrt3}{2}\right)^2
$$
--- разложение на множители на множестве $\mathbb C$.

$z=1$ --- корень кратности 4.

$z=\frac{-1\pm i\sqrt3}{2}$ --- корни кратности 2.

\end{document}
